% problems6.tex — Problem Set and Answers for Chapter 6

\begin{problemset}

\subsection*{Analytical Exercises}

\exercise[1] (Mean square limits and their uniqueness) \; In Chapter~2, we defined the mean
square convergence of a sequence of random variables $\{z_n\}$ to a random variable $z$.
In this and other questions, we will deal with sequences of random
variables with finite second moments. For those sequences, the definition of
mean square convergence is

\medskip
\noindent\textbf{Definition.} Let $\{z_n\}$ be a sequence of random variables with $\E(z_n^2) < \infty$.
We say that $\{z_n\}$ converges in mean square to a random variable $z$ with $\E(z^2) <
\infty$, written $z_n \to_{\mathrm{m.s.}} z$, if $\E[(z_n - z)^2] \to 0$ as $n \to \infty$.

\medskip
Take this result for granted in answering.

\begin{enumerate}[label=(\alph*)]
\item Show that, if $z_n \to_{\mathrm{m.s.}} z$ and $z_n \to_{\mathrm{m.s.}} z'$, then $\E[(z - z')^2] = 0$.
\textbf{Hint:} $z - z' = (z - z_n) + (z_n - z')$. Define $\|x\| = \sqrt{\E(x^2)}$. Then
\begin{align*}
  \|z - z'\|^2 &= \|z - z_n\|^2 + 2\,\E[(z - z_n)(z_n - z')] + \|z_n - z'\|^2 \\
  &\leq \|z - z_n\|^2 + 2\sqrt{\E[(z - z_n)^2]}\sqrt{\E[(z_n - z')^2]} + \|z_n - z'\|^2 \\
  &= (\|z - z_n\| + \|z_n - z'\|)^2.
\end{align*}
So $\|z - z'\| \leq \|z - z_n\| + \|z_n - z'\|$. This is called the \textbf{triangle inequality}.

\item Show that $\Pr(z = z') = 1$. \textbf{Hint:} Use Chebychev's inequality:
\[
  \Pr(|z - z'| > \varepsilon) \leq \frac{1}{\varepsilon^2}\,\E[\|z - z'\|^2].
\]
\end{enumerate}

\exercise[2] (Proof of Proposition~6.1) \; A well-known result about mean square convergence
(see, e.g., Brockwell and Davis, 1991, Proposition~2.7.1) is
\begin{enumerate}[label=(\roman*)]
\item $\{z_n\}$ converges in mean square if and only if $\E[(z_m - z_n)^2] \to 0$ as $m, n \to \infty$.
\item If $x_n \to_{\mathrm{m.s.}} x$ and $z_n \to_{\mathrm{m.s.}} z$, then
$\lim_{n \to \infty} \E(x_n) = \E(x)$ and $\lim_{n \to \infty} \E(x_n z_n) = \E(xz)$.
\end{enumerate}
Take this result for granted in answering.

\begin{enumerate}[label=(\alph*)]
\item Prove that \eqref{eq:6.1.6} converges in mean square as $n \to \infty$ under the hypothesis
of Proposition~6.1. \textbf{Hint:} Let $y_{t,n} = \mu + \sum_{j=0}^{n} \psi_j \varepsilon_{t-j}$.
What needs to be proved is that $\{y_{t,n}\}$ converges in mean square as $n \to \infty$
for each $t$. Given~(i) above, it is sufficient to prove that (assuming $m > n$)
\[
  \E\left[\left(\sum_{j=n+1}^{m} \psi_j \varepsilon_{t-j}\right)^2\right]
  = \sigma^2 \sum_{j=n+1}^{m} \psi_j^2 \to 0 \quad \text{as } m, n \to \infty.
\]
Use the fact that an absolutely summable sequence is \textbf{square summable}, i.e.,
$\sum |\psi_j| < \infty \;\Rightarrow\; \sum \psi_j^2 < \infty$,
and the fact that a sequence of real numbers $\{\alpha_n\}$ converges (to a finite
limit) if and only if $\alpha_n$ is a Cauchy sequence:
$\alpha_n \to \alpha \;\Leftrightarrow\; |\alpha_m - \alpha_n| \to 0$ as $m, n \to \infty$.
Set $\alpha_n = \sum_{j=0}^{n} \psi_j^2$.

\item Prove that $\E(y_t) = \mu$. \textbf{Hint:} Let $y_{t,n} = \mu + \sum_{j=0}^{n} \psi_j \varepsilon_{t-j}$ as in~(a). It was
shown in~(a) that $y_{t,n} \to_{\mathrm{m.s.}} y_t$.

\item (Proof of part~(b) of Proposition~6.1) \; Show that
\[
  \E[(y_t - \mu)(y_{t-j} - \mu)] = \lim_{n \to \infty} \E[(y_{t,n} - \mu)(y_{t-j,n} - \mu)].
\]
Have we shown that $\{y_t\}$ is covariance-stationary? [Answer: Yes.]

\item (Optional) \; Prove part~(c) of Proposition~6.1, taking the following facts
from calculus for granted.
\begin{enumerate}[label=(\roman*)]
\item If $\{a_j\}$ is absolutely summable, then $\{a_j\}$ is summable (i.e., $-\infty <
\sum_{j=0}^{\infty} a_j < \infty$) and $|\sum_{j=0}^{\infty} a_j| \leq \sum_{j=0}^{\infty} |a_j|$.
\item Consider a sequence with two subscripts, $\{a_{jk}\}$ ($j, k = 0, 1, 2, \ldots$).
Suppose $\sum_{j=0}^{\infty} |a_{jk}| < \infty$ for each $k$ and let $s_k \equiv \sum_{j=0}^{\infty} |a_{jk}|$. Suppose
$\{s_k\}$ is summable. Then
\[
  \left|\sum_{j=0}^{\infty}\left(\sum_{k=0}^{\infty} a_{jk}\right)\right| < \infty
  \quad \text{and} \quad
  \sum_{j=0}^{\infty}\left(\sum_{k=0}^{\infty} a_{jk}\right) = \sum_{k=0}^{\infty}\left(\sum_{j=0}^{\infty} a_{jk}\right) < \infty.
\]
\end{enumerate}
\textbf{Hint:} Derive
$\sum |\gamma_j| \leq \sigma^2 \sum_{j=0}^{\infty}\bigl(\sum_{k=0}^{\infty} |\psi_{j+k}| \cdot |\psi_k|\bigr) < \infty$.
\end{enumerate}

\exercise[3] (Autocovariances of $h(L)x_t$) \; We wish to derive the autocovariances of $\{y_t\}$
of Proposition~6.2. As in Proposition~6.2, assume throughout this question that
$\{x_t\}$ is covariance-stationary. First consider a weighted average of finitely many
terms: $y_{t,n} = h_0 x_t + h_1 x_{t-1} + \cdots + h_n x_{t-n}$.

\begin{enumerate}[label=(\alph*)]
\item Show that for this $\{y_{t,n}\}$,
\[
  \gamma_j = \sum_{k=0}^{n} \sum_{\ell=0}^{n} h_k\,h_\ell\,\gamma_{j-k+\ell}^x
\]
where $\gamma_j^x$ is the $j$-th order autocovariance of $\{x_t\}$.
Verify that this reduces to \eqref{eq:6.1.2a} if $\{x_t\}$ is white noise. \textbf{Hint:} If you find this
difficult, first set $n = 1$, and then increase $n$ to establish the pattern.

\item Now consider a weighted average of possibly infinitely many terms:
$y_t = h_0 x_t + h_1 x_{t-1} + \cdots$.
Taking Proposition~6.2(a) for granted, show that
\[
  \gamma_j = \sum_{k=0}^{\infty} \sum_{\ell=0}^{\infty} h_k\,h_\ell\,\gamma_{j-k+\ell}^x
  \quad \text{if } \{h_j\} \text{ is absolutely summable.}
\]
\textbf{Hint:} $y_{t,n} \to_{\mathrm{m.s.}} y_t$ as $n \to \infty$ by Proposition~6.2. Proceed as in~2(c).
\end{enumerate}

\exercise[4] (Homogeneous linear difference equations) \; Consider a $p$-th order homogeneous
linear difference equation:
\begin{equation}
  y_j - \phi_1 y_{j-1} - \phi_2 y_{j-2} - \cdots - \phi_p y_{j-p} = 0
  \quad (j = p, p+1, \ldots),
  \tag{1}
\end{equation}
where $(\phi_1, \phi_2, \ldots, \phi_p)$ are real constants with $\phi_p \neq 0$, with the initial condition
that $(y_0, y_1, \ldots, y_{p-1})$ are given. The associated polynomial equation is
\begin{equation}
  1 - \phi_1 z - \phi_2 z^2 - \cdots - \phi_{p-1} z^{p-1} - \phi_p z^p = 0.
  \tag{2}
\end{equation}
Let $\lambda_k$ ($k = 1, \ldots, K$) be the distinct roots of this equation (they can be
complex numbers) and $r_k$ be the multiplicity of $\lambda_k$. So $r_1 + r_2 + \cdots + r_K = p$.
The polynomial equation can be written as $\prod_{k=1}^{K} (1 - \lambda_k^{-1} z)^{r_k} = 0$.

As you verify below, the solution to the $p$-th order homogeneous equation
can be written as
\begin{equation}
  y_j = \sum_{k=1}^{K} \sum_{n=0}^{r_k - 1} c_{kn} \cdot (j)^n \lambda_k^{-j}.
  \tag{5}
\end{equation}
In the important special case where all the roots are distinct, $K = p$ and $r_k = 1$
for all $k$, and (5) becomes $y_j = \sum_{k=1}^{p} c_{k0}\,\lambda_k^{-j}$.

\begin{enumerate}[label=(\alph*)]
\item For the special case of $p = 2$ and distinct roots, (7) is
$y_j = c_{10}\,\lambda_1^{-j} + c_{20}\,\lambda_2^{-j}$.
Verify that this solves the homogeneous difference equation~(1) with $p = 2$
for any given $(c_{10}, c_{20})$. To determine $(c_{10}, c_{20})$, write down~(8) for $j = 0$
and $j = 1$, and solve for $c$'s as a function of $(y_0, y_1, \lambda_1, \lambda_2)$.

\item To learn (remember) how to deal with multiple roots, consider the example
$1 - z + \frac{1}{4}z^2 = 0$, whose root is~2 with multiplicity~2.
Verify that $y_j = c_{10} \cdot 2^{-j} + c_{11} \cdot j \cdot 2^{-j}$ solves the second-order homogeneous
difference equation $y_j - y_{j-1} + \frac{1}{4} y_{j-2} = 0$.

\item (Optional) \; Prove the following lemma.

\textbf{Lemma.} \emph{Let $0 \leq \xi < 1$ and let $n$ be a non-negative integer. Then there
exist two real numbers, $A$ and $b$, such that $\xi < b < 1$ and $(j)^n \xi^j < A b^j$
for $j = 0, 1, 2, \ldots$.}

\textbf{Hint:} Pick any $b$ such that $\xi < b < 1$. Because $b^j$ eventually gets
larger than $(j)^n \xi^j$ as $j$ increases, there exists some $j$, call it $J$, such that
$(j)^n \xi^j < b^j$ for all $j \geq J$.

\item Suppose the stability condition holds for the $p$-th order homogeneous linear
difference equation. Prove that there exist two numbers, $A$ and $b$, such
that $0 < b < 1$ and $|y_j| < A b^j$ for $j = 0, 1, 2, \ldots$.

\textbf{Hint:}
$|y_j| = \left|\sum_{k=1}^{K} \sum_{n=0}^{r_k-1} c_{kn} \cdot (j)^n \lambda_k^{-j}\right|
\leq \sum_{k=1}^{K} \sum_{n=0}^{r_k-1} |c_{kn}| \cdot (j)^n \cdot |\lambda_k^{-1}|^j$.
Since $|\lambda_k^{-1}| < 1$ by the stability condition, apply the lemma with $\xi = |\lambda_k^{-1}|$
and claim $(j)^n |\lambda_k^{-1}|^j < A_k (b_k)^j$ for some $A_k > 0$ and $|\lambda_k^{-1}| < b_k < 1$.
Set $A^* = \max\{A_k\}$ and $b = \max\{b_k\}$, so that $A_k(b_k)^j \leq A^* b^j$ for all $k$. Then define $A = cpA^*$.
\end{enumerate}

\exercise[5] (Yule--Walker equations) \; In the text, the autocovariances of an AR(1) process
were calculated from the $\text{MA}(\infty)$ representation. Here, we use the Yule--Walker
equations.

\begin{enumerate}[label=(\alph*)]
\item From $(6.2.1')$ on page~376, derive
\[
  \gamma_j - \phi \gamma_{j-1} = \E[\varepsilon_t \cdot (y_{t-j} - \mu)]
  \quad (j = 0, 1, 2, \ldots).
\]

\item (Easy) \; From the MA representation, show that
\[
  \E[\varepsilon_t \cdot (y_{t-j} - \mu)] = \begin{cases} \sigma^2 & (j = 0) \\ 0 & (j = 1, 2, \ldots). \end{cases}
\]

\item Derive the Yule--Walker equations:
\begin{align}
  \gamma_0 - \phi \gamma_1 &= \sigma^2, \tag{9} \\
  \gamma_j - \phi \gamma_{j-1} &= 0 \quad (j = 1, 2, \ldots). \tag{10}
\end{align}

\item Calculate $(\gamma_0, \gamma_1, \ldots)$ from (9) and (10). \textbf{Hint:} First set $j = 1$ in~(10) to
obtain $\gamma_1 = \phi \gamma_0$. This and~(9) can be solved for $(\gamma_0, \gamma_1)$.
\end{enumerate}

\exercise[6] (AR(1) without using filters) \; We wish to prove, without the help of Proposition~6.2,
that $\{y_t\}$ following the AR(1) equation~\eqref{eq:6.2.1}, with the stationarity
condition $|\phi| < 1$, must have an $\text{MA}(\infty)$ representation if it is covariance-stationary.

\begin{enumerate}[label=(\alph*)]
\item Let $x_t \equiv y_t - \mu$ so that $\E(x_t) = 0$. By successive substitution, derive from
$(6.2.1')$ on page~376
\[
  x_t = x_{t,n} + \phi^n x_{t-n}
\]
where $x_{t,n} = \varepsilon_t + \phi \varepsilon_{t-1} + \phi^2 \varepsilon_{t-2} + \cdots + \phi^{n-1} \varepsilon_{t-n+1}$.

\item Show that $x_{t,n} \to_{\mathrm{m.s.}} x_t$. \textbf{Hint:} What needs to be shown is that $\E[(x_{t,n} -
x_t)^2] \to 0$ as $n \to \infty$. $(x_{t,n} - x_t)^2 = \phi^{2n} x_{t-n}^2$. Because $\{y_t\}$ is assumed
to be covariance-stationary, $\E(x_{t-n}^2) < \infty$.

\item (Trivial) \; Recalling that the infinite sum $\sum_{j=0}^{\infty} \phi^j \varepsilon_{t-j}$
is \emph{defined} to be the mean square limit of the partial sum $x_{t,n}$, show that $y_t$
can be written as $y_t = \mu + \sum_{j=0}^{\infty} \phi^j \varepsilon_{t-j}$.
\end{enumerate}

\exercise[7] (Companion form) \; We wish to generalize the argument just developed in
question~6 to $\text{AR}(p)$. Without loss of generality, we can assume that $\mu = 0$ (or
redefine $y_t$ to be the deviation from the mean of $y_t$). Define
\[
  \underset{(p \times 1)}{\bm{\xi}_t} =
  \begin{bmatrix} y_t \\ y_{t-1} \\ y_{t-2} \\ \vdots \\ y_{t-p+1} \end{bmatrix}, \quad
  \underset{(p \times 1)}{\mathbf{v}_t} =
  \begin{bmatrix} \varepsilon_t \\ 0 \\ 0 \\ \vdots \\ 0 \end{bmatrix}, \quad
  \underset{(p \times p)}{\mathbf{F}} =
  \begin{bmatrix}
    \phi_1 & \phi_2 & \phi_3 & \cdots & \phi_{p-1} & \phi_p \\
    1 & 0 & 0 & \cdots & 0 & 0 \\
    0 & 1 & 0 & \cdots & 0 & 0 \\
    \vdots & \vdots & \vdots & & \vdots & \vdots \\
    0 & 0 & 0 & \cdots & 1 & 0
  \end{bmatrix}.
\]
Consider the first-order vector stochastic difference equation: $\bm{\xi}_t = \mathbf{F}\bm{\xi}_{t-1} + \mathbf{v}_t$.
In this system of $p$ equations, called the \textbf{companion form}, the first equation
is the same as the $\text{AR}(p)$ equation $(6.2.6')$ and the rest are just
identities about $y_{t-j}$ ($j = 1, 2, \ldots, p-1$).

\begin{enumerate}[label=(\alph*)]
\item By successive substitution, show that
$\bm{\xi}_t = \bm{\xi}_{t,n} + (\mathbf{F})^n \bm{\xi}_{t-n}$
where $\bm{\xi}_{t,n} \equiv \mathbf{v}_t + \mathbf{F}\mathbf{v}_{t-1} + \mathbf{F}^2 \mathbf{v}_{t-2} + \cdots + (\mathbf{F})^{n-1}\mathbf{v}_{t-n+1}$.

The \textbf{eigenvalues} or \textbf{characteristic roots} of a $p \times p$ matrix $\mathbf{F}$ are the roots
of the determinantal equation $|\mathbf{F} - \lambda \mathbf{I}_p| = 0$.
It can be shown that $|\mathbf{F} - \lambda \mathbf{I}_p| = \lambda^p - \phi_1 \lambda^{p-1} - \cdots - \phi_{p-1}\lambda - \phi_p$.
So the determinantal equation becomes the $p$-th order polynomial equation:
$\lambda^p - \phi_1 \lambda^{p-1} - \cdots - \phi_{p-1}\lambda - \phi_p = 0$.

\item Verify this for $p = 2$.

In the rest of this question, assume that all the roots of the $p$-th order polynomial
equation are distinct. Let $(\lambda_1, \lambda_2, \ldots, \lambda_p)$ be those distinct eigenvalues. We have from matrix algebra
that there exists a nonsingular $p \times p$ matrix $\mathbf{T}$ such that
$\mathbf{F} = \mathbf{T}\bm{\Lambda}\mathbf{T}^{-1}$, $\bm{\Lambda} = \diag(\lambda_1, \ldots, \lambda_p)$.

\item (Very easy) \; Prove that $(\mathbf{F})^n = \mathbf{T}(\bm{\Lambda})^n \mathbf{T}^{-1}$ for $n = 3$.

\item Show that $\bm{\xi}_{t,n} \to_{\mathrm{m.s.}} \bm{\xi}_t$ as $n \to \infty$ if the stationarity condition about the
$\text{AR}(p)$ equation is satisfied and if $\{y_t\}$ is covariance-stationary. \textbf{Hint:} The
stationarity condition requires that all the roots be less than~1 in absolute
value. By~(12) what needs to be shown is that $(\mathbf{F})^n \bm{\xi}_{t-n} \to_{\mathrm{m.s.}} \mathbf{0}$. Recall
that $\mathbf{z}_n \to_{\mathrm{m.s.}} \bm{\alpha}$ if the mean square convergence occurs component-wise.

\item Show that $y_t$ has the $\text{MA}(\infty)$ representation
$y_t = \psi_0 \varepsilon_t + \psi_1 \varepsilon_{t-1} + \psi_2 \varepsilon_{t-2} + \cdots$.
How is $\psi_j$ related to the characteristic roots? \textbf{Hint:} From~(d) it immediately
follows that $\bm{\xi}_t = \mathbf{v}_t + \mathbf{F}\mathbf{v}_{t-1} + (\mathbf{F})^2 \mathbf{v}_{t-2} + \cdots$.
The top equation of this is the $\text{MA}(\infty)$ representation for $y_t$.
\end{enumerate}

\exercise[8] (AR(1) that started at $t = 0$) \; In the text, the $\text{MA}(\infty)$ representation~\eqref{eq:6.2.2} of
an AR(1) process following $y_t = c + \phi y_{t-1} + \varepsilon_t$ assumes that the process is
defined for $t = -1, -2, \ldots$ as well as for $t = 0, 1, 2, \ldots$. Instead, consider
an AR(1) process that started at $t = 0$ with an initial value of $y_0$. Suppose that
$|\phi| < 1$ and that $y_0$ is uncorrelated with the subsequent white noise process
$(\varepsilon_1, \varepsilon_2, \ldots)$.

\begin{enumerate}[label=(\alph*)]
\item Write the variance of $y_t$ as a function of $\sigma^2$ ($\equiv \Var(\varepsilon_t)$), $\Var(y_0)$, and $\phi$.
\textbf{Hint:} By successive substitution,
\[
  y_t = (1 + \phi + \phi^2 + \cdots + \phi^{t-1})c + \varepsilon_t + \phi\varepsilon_{t-1} + \phi^2\varepsilon_{t-2}
  + \cdots + \phi^{t-1}\varepsilon_1 + \phi^t y_0.
\]
Let $\mu$ and $\{\gamma_j\}$ be the mean and autocovariances of the covariance-stationary
AR(1) process, so $\mu = c/(1-\phi)$ and $\{\gamma_j\}$ are as in \eqref{eq:6.2.5}.
Show that
\begin{gather*}
  \lim_{t \to \infty} \E(y_t) = \mu, \quad
  \lim_{t \to \infty} \Var(y_t) = \gamma_0, \quad \text{and} \quad
  \lim_{t \to \infty} \Cov(y_t, y_{t-j}) = \gamma_j.
\end{gather*}

\item Show that $\{y_t\}$ is covariance-stationary if $\E(y_0) = \mu$ and $\Var(y_0) = \gamma_0$.
\end{enumerate}

\exercise[9] (Proof of Proposition~6.8(a)) \; We wish to prove Proposition~6.8(a). By Chebychev's
inequality, it suffices to show that $\lim_{n \to \infty} \Var(\bar{y}) = 0$. From \eqref{eq:6.5.2},
\[
  \Var(\bar{y}) = \frac{\gamma_0}{n} + \frac{2}{n} \sum_{j=1}^{n} \Bigl(1 - \frac{j}{n}\Bigr) \gamma_j
  \leq \frac{\gamma_0}{n} + \frac{2}{n} \sum_{j=1}^{n} |\gamma_j|.
\]

\begin{enumerate}[label=(\alph*)]
\item (Optional) \; Prove the following result:
$\lim_{j \to \infty} a_j = 0 \;\Rightarrow\; \lim_{n \to \infty} \frac{1}{n} \sum_{j=1}^{n} a_j = 0$.

\textbf{Hint:} $\{a_j\}$ converges to~0. So (i)~$|a_j| < M$ for all $j$, and (ii)~for any given $\varepsilon >
0$ there exists a positive integer $N$ such that $|a_j| < \varepsilon/2$ for all $j > N$. So
$\frac{1}{n}\sum |a_j| < \frac{NM}{n} + \frac{\varepsilon}{2}$.
By taking $n$ large enough, $NM/n$ can be made less than $\varepsilon/2$.

\item Taking the result in~(a) as given, prove that $\Var(\bar{y}) \to 0$ as $n \to \infty$.
\end{enumerate}

\exercise[10] (Proof of Proposition~6.8(b))

\begin{enumerate}[label=(\alph*)]
\item (Optional) \; Prove the following result:
$\sum_{j=1}^{\infty} a_j < \infty \;\Rightarrow\; \lim_{n \to \infty} \frac{1}{n}\sum_{j=1}^{n} j\,a_j = 0$.

\textbf{Hint:}
$\sum_{j=1}^{n} j a_j = (a_1 + a_2 + \cdots + a_n) + (a_2 + a_3 + \cdots + a_n) + \cdots + (a_{n-1} + a_n) + a_n
= \sum_{j=1}^{n}\Big|\sum_{k=j}^{n} a_k\Big|$.

\item Taking the result in~(a) as given, prove Proposition~6.8(b).
\end{enumerate}

\subsection*{Empirical Exercises}

\exercise[1] Read at least pp.~115--119 (except the last two paragraphs of p.~119) of Bekaert
and Hodrick (1993) before answering. Data files DM.ASC (for the Deutsche
Mark), POUND.ASC (British Pound), and YEN.ASC (Japanese Yen) contain
weekly data on the following items:

\begin{quote}
Column 1: the date of the observation (e.g., ``19850104'' is January 4, 1985) \\
Column 2: the ask price of the dollar in units of the foreign currency in the
spot market on Friday of the current week ($S_t$) \\
Column 3: the ask price of the dollar in units of the foreign currency in the
30-day forward market on Friday of the current week ($F_t$) \\
Column 4: the bid price of the dollar in units of the foreign currency in the spot
market on the delivery date on a current forward contract ($S30_t$).
\end{quote}

The sample period is the first week of 1975 through the last week of 1989.
The sample size is~778. As in the text, define $s_t \equiv \log(S_t)$, $f_t \equiv \log(F_t)$,
$s30_t \equiv \log(S30_t)$. Define $\varepsilon_t \equiv s30_t - f_t$.

Pick your favorite currency to answer the following questions.

\begin{enumerate}[label=(\alph*)]
\item (Library/internet work) \; For the foreign currency of your choice, identify
the week when the absolute value of the forward premium is largest. For
that week, find some measure of the domestic one-month interest rate (e.g.,
the one-month CD rate) for the United States and the currency's home
country, to verify that the interest rate differential is as large as is indicated
in the forward premium.

\item (Correlogram of $\{\varepsilon_t\}$) \; Draw the sample correlogram of $\varepsilon_t$ with 40~lags.
Does the autocorrelation appear to vanish after 4~lags? (It is up to you to
decide whether to subtract the sample mean in the calculation of sample
correlations. Theory says the population mean is zero, which you might
want to impose in the calculation.)

\item (Is the log spot rate a random walk?) \; Draw the sample correlogram of
$s_{t+1} - s_t$ with 40~lags. For those 40 autocorrelations, use the Box--Ljung
statistic to test for serial correlation. Can you reject the hypothesis that $\{s_t\}$
is a random walk with drift?

\item (Unconditional test) \; Carry out the unconditional test. Can you replicate
the results of Table~6.1 for the currency?

\item (Optional, regression test with truncated kernel) \; Carry out the regression
test. Can you replicate the results of Table~6.2 for the currency?

\item (Bartlett kernel) \; Use the Bartlett kernel-based estimator of $\mathbf{S}$ to do the
regression test. Newey and West (1994) provide a data-dependent automatic
bandwidth selection procedure. Take for granted that the autocovariance
lag length determined by this procedure is 12 for yen/\$, 8 for DM/\$,
and 16 for Pound/\$. The standard error for the $f - s$ coefficient for yen/\$ should be 0.6815.

\item (Optional, VARHAC) \; Use the VARHAC estimator for $\hat{\mathbf{S}}$ in the regression
test. Set $p_{max}$ to the integer part of $n^{1/3}$. In Step~1 of the VARHAC, you
would estimate a bivariate VAR, and for each of the two VAR equations,
you would pick the lag length $p$ by BIC on the fixed sample of $t = p_{max} +
1, \ldots, T$. The information criterion could be either \eqref{eq:6.6.11} or, alternatively,
\[
  \log(SSR_p^{(k)}/(n - p_{max})) + pK \cdot \log(n - p_{max})/(n - p_{max}).
\]
For yen/\$, the lag length selected should be 4 for the first VAR equation and 6 for the second.
The standard error for the $f - s$ coefficient should be 0.8027.
\end{enumerate}

\exercise[2] (Continuation of Fama exercise of Chapter~2) \; Now we know how to use
the three-month T-bill rate to test the efficiency of the U.S.\ Treasury bills market
using monthly data. Let

$\pi_t =$ inflation rate (in percent, annual rate) over the three-month
period ending in month $t - 1$,

$R_t =$ three-month T-bill rate (in percent, annual rate) at the beginning
of month $t$.

Consider the Fama regression for testing market efficiency:
\[
  \pi_{t+3} = \beta_0 + \beta_1 R_t + \varepsilon_t.
\]

\begin{enumerate}[label=(\alph*)]
\item (Trivial) \; Why is the dependent variable $\pi_{t+3}$ rather than, say, $\pi_t$?

\item Show that, under the efficient market hypothesis, $\beta_1 = 1$ and the autocovariance
of $\{\varepsilon_t\}$ vanishes after the second lag (i.e., $\gamma_j = 0$ for $j > 2$).

\item Let $r_{t+3} \equiv R_t - \pi_{t+3}$ be the ex-post real rate on the three-month T-bill.
Use \emph{PAI3} in MISHKIN.ASC as the measure of the three-month inflation
rate to calculate the ex-post real rate. Draw the sample correlogram for the sample period of 1/53--7/71. What is
the prediction of the efficient market hypothesis about the correlogram?

\item Test market efficiency by estimating the Fama regression for the sample
period of 1/53--7/71. (The sample size should be~223; do not throw away
two-thirds of the observations, as Fama did in his Tables~6--8.) Use \eqref{eq:6.6.4}
with $q = 2$ for $\hat{\mathbf{S}}$.
\end{enumerate}

\end{problemset}

\begin{answers}

\subsection*{Analytical Exercises}

\answer{2d} Since $\{\psi_j\}$ is absolutely summable, $\psi_j \to 0$ as $j \to \infty$. So for any $j$, there
exists an $A > 0$ such that $|\psi_{j+k}| \leq A$ for all $j, k$. So $|\psi_{j+k} \cdot \psi_k| \leq A|\psi_k|$.
Since $\{\psi_k\}$ (and hence $\{A\psi_k\}$) is absolutely summable, so is $\{\psi_{j+k} \cdot \psi_k\}$ ($k =
0, 1, 2, \ldots$) for any given $j$. Thus by~(i),
\[
  |\gamma_j| = \sigma^2 \left|\sum_{k=0}^{\infty} \psi_{j+k}\,\psi_k\right|
  \leq \sigma^2 \sum_{k=0}^{\infty} |\psi_{j+k}\,\psi_k|
  = \sigma^2 \sum_{k=0}^{\infty} |\psi_{j+k}|\,|\psi_k| < \infty.
\]
Now set $a_{jk}$ to $|\psi_{j+k}| \cdot |\psi_k|$. Then
\[
  \sum_{j=0}^{\infty} |a_{jk}| = \sum_{j=0}^{\infty} |\psi_{j+k}|\,|\psi_k|
  \leq |\psi_k| \sum_{j=0}^{\infty} |\psi_j| < \infty.
\]
Let $M \equiv \sum |\psi_j|$ and $s_k \equiv |\psi_k| \sum_{j=0}^{\infty} |\psi_{j+k}|$.
Then $|s_k| \leq |\psi_k| \cdot M$ and $\{\psi_k\}$ is absolutely summable.
Therefore, by~(ii),
$\sum_{j=0}^{\infty}\bigl(\sum_{k=0}^{\infty} |\psi_{j+k}| \cdot |\psi_k|\bigr) < \infty$.
This times $\sigma^2$ is $\sum_{j=0}^{\infty} |\gamma_j|$. So $\{\gamma_j\}$ is absolutely summable.

\answer{7e} $\psi_j$ is the $(1, 1)$ element of $\mathbf{T}(\bm{\Lambda})^j \mathbf{T}^{-1}$.

\answer{8a}
\begin{align*}
  \E(y_t) &= \frac{1 - \phi^t}{1 - \phi}\,c + \phi^t\,\E(y_0)
  = (1 - \phi^t)\mu + \phi^t\,\E(y_0)
  = \mu + \phi^t \cdot [\E(y_0) - \mu], \\
  \Var(y_t) &= (1 + \phi^2 + \phi^4 + \cdots + \phi^{2t-2})\sigma^2 + \phi^{2t}\,\Var(y_0)
  \quad (\text{since } \Cov(\varepsilon_t, y_0) = 0) \\
  &= \frac{1 - \phi^{2t}}{1 - \phi^2}\,\sigma^2 + \phi^{2t}\,\Var(y_0)
  = (1 - \phi^{2t})\gamma_0 + \phi^{2t}\,\Var(y_0) \\
  &= \gamma_0 + \phi^{2t} \cdot [\Var(y_0) - \gamma_0], \\
  \Cov(y_t, y_{t-j}) &= (\phi^j + \phi^{j+2} + \cdots + \phi^{2t-j-2})\sigma^2 + \phi^{2t-j}\,\Var(y_0) \\
  &= \phi^j \cdot \frac{1 - \phi^{2(t-j)}}{1 - \phi^2}\,\sigma^2 + \phi^{2t-j}\,\Var(y_0) \\
  &= \gamma_0 \phi^j + \phi^{2t-j} \cdot [\Var(y_0) - \gamma_0]
  = \gamma_j + \phi^{2t-j} \cdot [\Var(y_0) - \gamma_0].
\end{align*}

\subsection*{Empirical Exercises}

\answer{1c} The Ljung--Box statistic with 40~lags is 84.0 with a $p$-value of 0.0058\%. So the
hypothesis that the spot rate is a random walk with drift can be rejected. See
Figure~6.6 for the correlogram for yen/\$.

\begin{figure}[H]
\centering
\includegraphics[width=0.9\textwidth,trim=50 350 50 350,clip]{figures/fig_6_6.png}
\caption{Correlogram of $s_{t+1} - s_t$, Yen/Dollar}
\label{fig:6.6}
\end{figure}

\answer{1f} For the yen/\$ exchange rate, the standard errors are 3.72 for the constant and
0.68 for the $f - s$ coefficient. The Wald statistic for the hypothesis that $\beta_0 = 0$
and $\beta_1 = 1$ is 21.85 ($p = 0.002\%$).

\answer{2c} The correlogram for the three-month ex-post real rate is shown in Figure~6.7.
Theory predicts that serial correlation vanishes after the second lag. The
correlogram is mostly consistent with this prediction.

\begin{figure}[H]
\centering
\includegraphics[width=0.9\textwidth,trim=50 350 50 350,clip]{figures/fig_6_7.png}
\caption{Correlogram of Three-Month Ex-Post Real Rate}
\label{fig:6.7}
\end{figure}

\end{answers}
